% !TeX root = ../main.tex
\chapter{Background and Related Work}
\label{chap:background}

\textbf{Draft outline.} Discuss continuous delivery, regression testing,
AI-enabled systems, and the literature that motivates CD4AI.

% TODO: Define foundation models here, using Bommasani:FoundationModels:2021.
% Distinguish foundation models from pretrained models in general, and explain
% the role of language foundation models in the applications discussed later.
% Once this explanation is written, replace the definition currently opening
% Section 4.1 (sec:cd4ai-context) with a reference to this chapter, keeping the
% application-development context in Chapter 4.

% Temporary citation example; replace this paragraph when drafting the chapter.
The existing work-plan bibliography includes a book on continuous
delivery~\cite{Humble:ContinuousDelivery:2010}.

\section{Testing and Evaluation Terminology}
\label{sec:background-evaluation}

The terms \emph{test} and \emph{evaluation} overlap in AI engineering practice.
For example, \textcite{Grace:DemystifyingEvals:2026} describe an evaluation,
often abbreviated to \emph{eval}, as a test that supplies an input to an AI
system and assesses its output using grading logic. The terminology below
establishes a convention for this thesis rather than a universal distinction
between the two terms.

\subsection{Evaluations, Cases, and Tests}

\begin{description}
  \item[Evaluation.] The assessment of application behavior against specified
    criteria. Results may include scores, pass/fail judgments, and summaries
    across cases or repeated executions.
  \item[Evaluation case.] A retained input or scenario, the context needed to
    exercise it, and criteria for assessing the resulting behavior. An expected
    outcome may be expressed through assertions or a rubric rather than a
    single reference answer.
  \item[Test.] An executable check of expected behavior that may be used within
    an evaluation. Calling a check a test does not imply that the application
    or the assessment method is deterministic.
  \item[Evaluation suite.] A collection of cases used to assess selected
    behaviors or capabilities. A dataset may store cases, while executing an
    evaluation also requires an application, an execution procedure, and a
    means of assessing the results.
\end{description}

This thesis uses \emph{evaluation} as the general term for behavioral assessment
and retains \emph{testing} when discussing executable checks and software
testing practices. The choice of term does not by itself determine whether a
result should block deployment.

\subsection{Confirmation, Capability, and Regression Evaluation}
\label{sec:background-evaluation-purpose}

Confirmation testing assesses whether a proposed fix resolves an identified
defect. Regression testing checks for adverse effects of changes on existing
behavior~\cite[Section~2.2.3]{ISTQB:FoundationSyllabus:2024}. These purposes
should be distinguished from the time at which a case is written: a test
intended for regression protection may be created before the defect is fixed,
and a newly written test may expose a regression that has already occurred.

In the context of agents, \textcite{Grace:DemystifyingEvals:2026} distinguish
capability evaluations, which assess abilities that may still need improvement,
from regression evaluations, which check whether established capabilities are
preserved. They also describe how cases initially used to guide improvement
can later support regression evaluation as the system becomes capable of
handling them.

Accordingly, a case documenting an unresolved defect remains useful even if it
currently fails. It can record the desired behavior and help assess proposed
fixes. Once that behavior is established, the same case can help detect its
deterioration under later changes. A continuing failure on a known unresolved
defect does not, by itself, demonstrate a new regression.

\subsection{Retained Cases and Deployment Decisions}

For CD4AI, retaining an evaluation case and deciding whether its result blocks
deployment are separate decisions. Cases for unresolved defects should retain
their expected behavior and make the unresolved status visible. Depending on
the application's acceptance policy, such a failure may block a release or be
reported as a known limitation whose resolution has been deferred. Retaining
the case does not mean that the defect has been fixed or that the observed
incorrect behavior has become acceptable.

In the title of this thesis, \emph{continuous evaluation} refers to assessment
repeatedly integrated into the delivery and feedback cycle. It does not require
uninterrupted execution or a particular evaluation framework. The cycle
supports evaluation of improvements as well as detection of regressions.
